% piTantum -- Appendici tecniche (Parte II)
% Per sviluppatori, manutentori, integratori.
% ASCII puro, accenti via macro \`e \'e \`a.
% Manuale tecnico del progetto piTantum (Tempus Tantum;
% codename interno: timetable).
% ASCII puro, accenti via macro \`e \'e \`a, virgolette ``...''.
\documentclass[a4paper,11pt,oneside]{report}

\usepackage[T1]{fontenc}
\usepackage[utf8]{inputenc}
\usepackage[italian]{babel}
\usepackage{lmodern}
\usepackage{microtype}
\usepackage{geometry}
\geometry{margin=2.4cm,headheight=14pt}
\usepackage{xcolor}
\usepackage{graphicx}
\usepackage{tikz}
\usetikzlibrary{positioning,arrows.meta,shapes.geometric,fit,calc}
\usepackage{listings}
\usepackage{enumitem}
\usepackage{array}
\usepackage{longtable}
\usepackage{booktabs}
\usepackage{fancyhdr}
\usepackage{titlesec}
\usepackage[hidelinks,bookmarksopen=true]{hyperref}

\definecolor{cHard}{RGB}{220,38,38}
\definecolor{cSoft}{RGB}{217,119,6}
\definecolor{cPref}{RGB}{37,99,235}
\definecolor{cEnf}{RGB}{6,95,70}
\definecolor{cAllow}{RGB}{16,185,129}
\definecolor{cBox}{RGB}{229,231,235}
\definecolor{cCode}{RGB}{30,41,59}

\lstset{
  basicstyle=\ttfamily\footnotesize,
  keywordstyle=\color{cPref}\bfseries,
  commentstyle=\color{gray}\itshape,
  stringstyle=\color{cEnf},
  showstringspaces=false,
  breaklines=true,
  breakatwhitespace=true,
  columns=fullflexible,
  frame=single,
  rulecolor=\color{cBox},
  framesep=4pt,
  xleftmargin=8pt,
  xrightmargin=4pt,
  numbers=none,
}
\lstdefinestyle{json}{language=,morekeywords={true,false,null}}
\lstdefinestyle{py}{language=Python}
\lstdefinestyle{shell}{language=bash,morekeywords={apt,brew,dnf,pacman,pip,npm,git,curl}}

\titleformat{\chapter}[hang]{\Huge\bfseries}{\thechapter.}{12pt}{}
\titleformat{\section}{\Large\bfseries}{\thesection}{10pt}{}
\titleformat{\subsection}{\large\bfseries}{\thesubsection}{8pt}{}

\pagestyle{fancy}
\fancyhf{}
\fancyhead[L]{\textit{piTantum -- manuale tecnico}}
\fancyhead[R]{\thepage}
\renewcommand{\headrulewidth}{0.4pt}

\title{}
\author{}
\date{}



\begin{document}

% =========== Cover (Parte II) ===========
\begin{titlepage}
\thispagestyle{empty}
\centering
\vspace*{1.4cm}

{\Huge\bfseries $\pi$\,Tantum\par}
\vspace{0.2cm}
{\Large alias \emph{Tempus Tantum}\par}
\vspace{0.4cm}
{\Large\itshape sistema di generazione orario scolastico\par}

\vspace{1.0cm}
\rule{0.5\textwidth}{0.4pt}

\vspace{1.4cm}
{\LARGE Appendici tecniche\par}
\vspace{0.4cm}
{\Large\itshape Parte II --- per sviluppatori\par}

\vspace{2cm}

\begin{minipage}{0.8\textwidth}
\centering\small
Riferimento tecnico completo: architettura del sistema,
modello dati con relazioni e vincoli FK, API REST, grammatica
formale del DSL e traduzione AST $\to$ CP-SAT, strategie di
ottimizzazione (CP-SAT, decomposizione spettrale,
metaeuristiche), diagnostica statistica, guida all'estensione
del sistema.
\end{minipage}

\vspace{1.4cm}
\begin{tikzpicture}
  \node[draw=cBox, very thick, rounded corners=4pt, minimum width=3cm,
        minimum height=1cm, fill=cBox!30] (a) at (0,0) {\textbf{CP-SAT}};
  \node[draw=cBox, very thick, rounded corners=4pt, minimum width=3cm,
        minimum height=1cm, fill=cBox!30, right=0.6cm of a] (b)
        {\textbf{spectral}};
  \node[draw=cBox, very thick, rounded corners=4pt, minimum width=3cm,
        minimum height=1cm, fill=cBox!30, right=0.6cm of b] (c)
        {\textbf{LNS+SA+TS+ILS}};
  \draw[-{Stealth[length=3mm]}, thick] (a) -- (b);
  \draw[-{Stealth[length=3mm]}, thick] (b) -- (c);
\end{tikzpicture}

\vfill

{\large Giovanni Borghi\par}
\vspace{0.2cm}
{\large \today\par}
\end{titlepage}

\tableofcontents

% =========== Appendici tecniche ===========
\chapter{Architettura}

\section{Stack}

\begin{itemize}[leftmargin=*]
  \item \textbf{Backend}: Python 3.11+, FastAPI, SQLAlchemy 2.0,
        Pydantic 2, uvicorn, OR-Tools, NumPy, scikit-learn,
        Faker.
  \item \textbf{Frontend}: SvelteKit (Svelte 4), Vite 5, Tailwind CSS,
        JavaScript puro.
  \item \textbf{DB}: SQLite via SQLAlchemy. Migrazioni idempotenti
        nel codice (no Alembic).
  \item \textbf{Engine I/O}: pickle Python; \texttt{engine\_io.py}
        converte fra DB e dict per il solver.
\end{itemize}

\section{Diagramma a blocchi}

\begin{figure}[h]
\centering
\begin{tikzpicture}[node distance=0.6cm and 0.6cm,
  every node/.style={draw, rounded corners=4pt, very thick,
                     minimum width=4.5cm, minimum height=1cm,
                     align=center, font=\small}]
  \node[fill=cBox!30] (fe) {SvelteKit / Vite\\\textit{16 tabs}};
  \node[fill=cBox!30, below=of fe] (be) {FastAPI backend\\\textit{18 router, 118 endpoint}};
  \node[fill=cBox!30, below=of be] (db) {SQLite\\\textit{$\sim$30 tabelle}};
  \node[fill=cBox!30, below=of db] (en) {CP-SAT engine\\\textit{decomp. spettrale + meta}};
  \draw[-{Stealth[length=3mm]}, thick] (fe) -- node[draw=none, right]
        {\footnotesize \texttt{/api/*} via Vite proxy} (be);
  \draw[-{Stealth[length=3mm]}, thick] (be) -- node[draw=none, right]
        {\footnotesize SQLAlchemy 2.0} (db);
  \draw[-{Stealth[length=3mm]}, thick] (be) to[bend left=45]
        node[draw=none, right]
        {\footnotesize \texttt{engine\_io.py} (pickle)} (en);
\end{tikzpicture}
\caption{Architettura a blocchi della webapp.}
\label{fig:arch}
\end{figure}

\section{Struttura cartelle}

\begin{lstlisting}[style=shell]
timetable/
+-- README.md                    -- landing page
+-- docs/                        -- manuale tecnico
+-- experiments/                 -- solver code + pickle snapshots
|   +-- cpsat_v2_assignment.py   -- Phase A
|   +-- cpsat_v2_timetable.py    -- Phase B
|   +-- decomposition_spectral_v2.py
|   +-- metaheuristics.py        -- LNS / SA / TS / ILS
|   +-- classroom_assignment.py  -- Phase 4
+-- schedule/                    -- legacy single-file prototypes
+-- webui/
|   +-- start.bat                -- launcher Windows
|   +-- start.sh                 -- launcher Linux/macOS
|   +-- stop.sh                  -- stop helper
|   +-- backend/                 -- FastAPI app
|   +-- frontend/                -- SvelteKit
|   +-- data/timetable.db        -- SQLite
|   +-- docs/                    -- guide brevi user-facing
+-- proposals/                   -- design notes + benchmarks
+-- *.pdf                        -- reference papers
\end{lstlisting}

\section{Lifecycle del backend}

\texttt{webui/backend/main.py} definisce un \texttt{lifespan} che
chiama \texttt{init\_db()} allo startup. Esso fa:

\begin{enumerate}
  \item \texttt{Base.metadata.create\_all(bind=engine)} -- crea le
        tabelle mancanti.
  \item \texttt{\_apply\_lightweight\_migrations()} -- ALTER TABLE
        idempotenti per i campi aggiunti in versioni successive
        (esempio: \texttt{school\_classes.curriculum\_id},
        \texttt{teachers.last\_name/first\_name/nickname}, ecc.).
        Ogni ALTER \`e guardato da \texttt{PRAGMA table\_info} per non
        duplicare.
\end{enumerate}

\section{Architettura del solver: dal DSL a CP-SAT}
\label{sec:arch-dsl-cpsat}

Step 1--4 del piano multi-day hanno introdotto un'architettura
\emph{a strati uniforme} per tutti gli usi del solver. Il
diagramma sotto riassume i flussi di trasformazione.

\begin{verbatim}
+-----------------------------------------------------------+
| UI (SvelteKit)                                            |
|   - Tab Vincoli: matrici + DSL + Plessi                   |
|   - Tab Workflow: pipeline (mode, granularity)            |
+-----------------+----------------------------+------------+
                  |                              |
        POST /api/constraints/general     POST /api/optimize/*
                  |                              |
                  v                              v
+-----------------+----------------+   +---------+---------+
| dsl_parser (general_dsl.py)      |   | router            |
|   parse() -> AST                  |   | run_async         |
+-----------------+----------------+   +---------+---------+
                  |                              |
                  v                              v
+-----------------+----------------+   +---------+---------+
| dsl_translator                   |   | engine pipeline   |
|   tabelle special-purpose        |   | (mono / decomp /  |
|   -> stringhe DSL canoniche       |   |  CG / BP)         |
+-----------------+----------------+   +---------+---------+
                  |                              |
                  +------------- merged ---------+
                                  |
                                  v
+--------------------------------+--------------------------+
| dsl_to_cpsat.DSLConstraintCompiler                         |
|   AST + slot dict + plessi_data                            |
|   -> add_assertion() / add_implication() su CpModel        |
+--------------------------------+--------------------------+
                                 |
              +------------------+------------------+
              v                  v                  v
       MonolithicSolver   PhaseBDaySolver   BP Pricer (9 grani)
       (cpsat_v2_*) +     (legacy +         + RF nodi
       seed via DSL       DSL augment)
              |                  |                  |
              +------------------+------------------+
                                 |
                                 v
                         soluzione (Lessons)
                                 |
                                 v
                          DSL evaluator post-hoc
                          (HARD: feasibility check)
                          (SOFT: contributo allo
                                  score globale)
\end{verbatim}

I punti chiave:
\begin{itemize}
\item \textbf{Un parser, un AST}. Tutte le sorgenti (matrici,
  DSL UI, regole plesso, vincoli legacy seedati) vengono
  unificate in stringhe DSL canoniche e parsate da
  \texttt{general\_dsl.py}.
\item \textbf{Un compilatore}. Il
  \texttt{DSLConstraintCompiler} \`e la porta unica verso
  CP-SAT. Mono solver, pricer di BP, PhaseBDaySolver, nodi
  Ryan-Foster: tutti applicano lo stesso compilatore al loro
  modello, garantendo coerenza dei vincoli fra le diverse
  pipeline.
\item \textbf{Zero-drift}.
  \texttt{seed\_implicit\_hardcoded(profs)} traduce in DSL i
  vincoli HARD legacy hardcoded. Le tuple di slot proibite
  emesse via DSL coincidono bit-per-bit con quelle del path
  legacy: la migrazione \`e progressiva e reversibile.
\item \textbf{SOFT post-hoc, HARD up-front}. I HARD del DSL
  diventano vincoli CP-SAT prima della solve. I SOFT del DSL
  sono attualmente valutati \emph{post-hoc}: il loro costo
  contribuisce allo score globale ma non orienta ancora la
  ricerca CP-SAT (la wired-up del SOFT in obiettivo \`e
  pianificata nei prossimi step).
\end{itemize}

% =========== Cap 3: Modello dati ===========

\chapter{Modello dati}

Tutto vive in \texttt{webui/backend/models.py}. Le tabelle si
raggruppano in cinque famiglie tematiche.

\section{Famiglie di tabelle}

\subsection*{Anagrafiche}
\texttt{subjects}, \texttt{teachers}, \texttt{school\_classes},
\texttt{classrooms}, \texttt{curricula}, \texttt{students},
\texttt{study\_groups}.

\subsection*{Tag (M-N)}
\texttt{classroom\_tags} + \texttt{classroom\_tag\_assignments},
\texttt{student\_tags} + \texttt{student\_tag\_assignments}. Nome
unique lower-cased, descrizione opzionale, cancellazione cascata.

\subsection*{Assegnazione}
\texttt{class\_subjects}, \texttt{curriculum\_subject\_hours},
\texttt{teacher\_subjects}, \texttt{subject\_group\_weights},
\texttt{assignments}.

\subsection*{Disponibilit\`a / vincoli}
\texttt{teacher\_unavailability}, \texttt{class\_unavailability},
\texttt{classroom\_unavailability},
\texttt{logical\_unavailabilities},
\texttt{curriculum\_logical\_constraints},
\texttt{classroom\_subject\_preferences},
\texttt{teacher\_classroom\_preferences},
\texttt{coteaching\_rules}.

\subsection*{Soluzioni e scheduling}
\texttt{solutions}, \texttt{lessons}, \texttt{day\_counts}.

\subsection*{Coverage e Run}
\texttt{absences}, \texttt{substitute\_assignments},
\texttt{runs}, \texttt{run\_logs}, \texttt{app\_state}.

\subsection*{Viste salvate}
\texttt{saved\_views}: query DSL + sort multi-livello salvati per
una entit\`a lista, esposti dal dropdown "Viste salvate" del
componente \emph{SortableQueryableList}. UNIQUE \texttt{(entity,
name)}.

\section{Diagramma relazioni}

\begin{figure}[h]
\centering
\begin{tikzpicture}[node distance=0.5cm and 0.4cm,
  rect/.style={draw, rounded corners=2pt, thick, fill=cBox!20,
               minimum width=2.4cm, minimum height=0.6cm,
               align=center, font=\footnotesize}]
  \node[rect] (sub) {subjects};
  \node[rect, right=of sub] (ts) {teacher\_subjects};
  \node[rect, right=of ts] (te) {teachers};
  \node[rect, below=of te] (asn) {assignments};
  \node[rect, left=of asn] (sc) {school\_classes};
  \node[rect, left=of sc] (cur) {curricula};
  \node[rect, below=of asn] (sol) {solutions};
  \node[rect, below=of sol] (les) {lessons};
  \node[rect, below=of cur] (cs) {class\_subjects};
  \node[rect, below=of sc] (cu) {class\_unavail.};
  \node[rect, right=of les] (rm) {classrooms};
  \node[rect, above=of rm] (cp) {classroom\_subj\_pref};
  \draw (sub) -- (ts) -- (te);
  \draw (te) -- (asn);
  \draw (sc) -- (asn);
  \draw (sc) -- (cs);
  \draw (sc) -- (cu);
  \draw (cur) -- (sc);
  \draw (asn) -- (sol);
  \draw (sol) -- (les);
  \draw (rm) -- (cp);
  \draw (sub) to[out=270,in=180,looseness=1.4] (cp);
\end{tikzpicture}
\caption{Relazioni rilevanti del modello dati (semplificato).}
\label{fig:datamodel}
\end{figure}

\section{Migrazioni idempotenti}

Pattern (in \texttt{db.py}):

\begin{lstlisting}[style=py]
if insp.has_table("school_classes") \
        and not has_column("school_classes", "curriculum_id"):
    conn.execute(text(
        "ALTER TABLE school_classes ADD COLUMN curriculum_id INTEGER"
    ))
\end{lstlisting}

Ogni ALTER \`e guardato da una check su \texttt{PRAGMA table\_info}.
Quando aggiungi una colonna nuova, segui lo stesso pattern: il DB
preesistente sopravvive senza interventi manuali.

\section{Pickle dell'engine}

\texttt{engine\_io.py} converte fra DB e tre forme pickle:

\begin{lstlisting}[style=py]
school = {
  'profile':   str,
  'classes':   [{name, year, section, curriculum, subjects:{subj:ore}}],
  'teachers':  [{name, group, max_hours, free_day,
                 weights:{subj:int}}],
  'cconcorsopersubject': {subj: {group: weight}},
  'curriculum_scores':   {curriculum: int},
  'curricula':           [...],
  'students':            [...],
  'groups':              [...],
}
profs = {
  teacher_name: {
    'classi':  {class_name: {subject: {'ore': N}}},
    'glibero': [d1, d2, d3]
  }
}
solution = {(prof, class, subj, day, hour): 0|1}
\end{lstlisting}

% =========== Cap 4: Vincoli ===========

\chapter{DSL generico per i vincoli}
\label{ch:dsl}
\label{ch:dsl_generico_per_i_vincoli}

Linguaggio compatto, dichiarativo, totalmente componibile per
esprimere vincoli HARD/SOFT/PREFERRED/ENFORCED su qualunque
combinazione logica di predicati sul modello dei dati.
Implementato in \texttt{webui/backend/utils/general\_dsl.py}
(parser ricorsivo discendente + AST + evaluator post-hoc, $\sim$250
LOC, no dipendenze esterne). Esposto via
\texttt{POST /api/constraints/general}, validato live da
\texttt{POST /api/constraints/general/validate}, accessibile dal
modal ``+ Nuovo vincolo DSL'' del tab Vincoli.

\section{Filosofia}

Prima del DSL generico il sistema aveva quattro sub-DSL
specializzati (query, logical, objective, introspettivo). Stessa
grammatica implementata 4 volte. Il DSL generico unifica tutto
sotto una sola grammatica con quattro ``sapori'' di compilazione
(LIST, LOGIC, OBJECTIVE, EVAL): \emph{un parser, molti
compilatori}. La sintassi di
\texttt{forall x in S where Filter: Body} e dei predicati atomici
(\texttt{==}, \texttt{!=}, \texttt{<}, \texttt{in}, \dots) e'
identica in tutti i contesti; cambia solo il backend di
traduzione.

Il DSL e' interpretato \emph{post-hoc}: parsea l'espressione e
la valuta DOPO che la pipeline (Phase A, Phase B, metaeuristiche)
ha prodotto una soluzione candidata. Le violazioni HARD sono
rilevate dal Feasibility Check; le SOFT alimentano lo score
globale.

\section{Grammatica BNF}

\begin{lstlisting}[basicstyle=\ttfamily\small]
expr           ::= iff_expr
iff_expr       ::= implies_expr ( "<=>" implies_expr )*
implies_expr   ::= or_expr     ( "=>"  or_expr )*
or_expr        ::= and_expr    ( ("OR"|"or") and_expr )*
and_expr       ::= not_expr    ( ("AND"|"and") not_expr )*
not_expr       ::= ("NOT"|"not") not_expr | atom

atom           ::= quant_expr | count_expr | comparison
                 | call | bool_lit | "(" expr ")"

quant_expr     ::= ("forall"|"exists") IDENT "in" source
                   ["where" or_expr] ":" expr

count_expr     ::= "count" IDENT "in" source ["where" or_expr]
                   ( ":" comparison | op value )

source         ::= IDENT                  -- literal name
                 | IDENT ("." IDENT)+     -- path (runtime)

comparison     ::= value ( op value
                         | "in" "[" value ("," value)* "]"
                         | "in" path
                         | "not_in" "[" value ("," value)* "]"
                         | "not_in" path )

op             ::= "==" | "!=" | "<" | "<=" | ">" | ">="
value          ::= IDENT ( "." IDENT )* | NUM | STRING
                 | BOOL | function_call
function_call  ::= IDENT "(" [arg_list] ")"
arg_list       ::= arg ("," arg)*
arg            ::= IDENT "==" value | value
\end{lstlisting}

Precedenze (dal pi\`u basso al pi\`u alto):
\texttt{<=>}, \texttt{=>}, \texttt{OR}, \texttt{AND}, \texttt{NOT},
confronti, \texttt{.} (member access). Le parentesi sovrascrivono.

\section{Sorgenti iterabili}

\begin{tabular}{ll}
\toprule
\textbf{Sorgente} & \textbf{Tipo elemento} \\
\midrule
\texttt{lessons}     & lezioni della soluzione attiva \\
\texttt{assignments} & cattedre (Assignment rows) \\
\texttt{teachers}    & docenti (\texttt{name, group, max\_hours, subject}) \\
\texttt{classes}     & classi (\texttt{name, year, curriculum, n\_students}) \\
\texttt{classrooms}  & aule (\texttt{name, kind, type, tags}) \\
\texttt{students}    & studenti (con \texttt{tags[]} e \texttt{groups[]}) \\
\texttt{subjects}    & materie \\
\texttt{curricula}   & indirizzi di studio \\
\texttt{groups}      & gruppi articolati \\
\texttt{slots}       & le 36 \texttt{(day, hour)} della settimana \\
\texttt{days}        & i 6 giorni (Lun..Sab; index 1..6) \\
\texttt{hours}       & le 6 ore (8..13) \\
\bottomrule
\end{tabular}

\textbf{Sorgente-path}: dalla v0.5 puoi scrivere
\texttt{exists g in s.groups: g.name == "X"} dove \texttt{s.groups}
e' un attributo che restituisce una lista; il parser accetta
qualunque path puntato dopo \texttt{in}.

\section{Galleria di esempi}

Ogni esempio sotto e' stato validato live contro il parser
corrente; copia-incolla nell'editor del modal ``+ Nuovo vincolo
DSL'' e funziona.

\subsection{Vincoli docente (10 esempi)}

\begin{enumerate}
\item Borghi non insegna mai in 1A alla 6a ora:
\begin{lstlisting}
forall l in lessons where l.teacher == Borghi
                       and l.class == 1A:
    l.hour != 13
\end{lstlisting}

\item Borghi: max 4 ore/giorno totali:
\begin{lstlisting}
forall d in days:
    count l in lessons where l.teacher == Borghi
                          and l.day == d.index: l <= 4
\end{lstlisting}

\item Lab fisica: ogni docente di Fisica esattamente 1 ora a
settimana in lab\_fisica:
\begin{lstlisting}
forall t in teachers where t.subject == "Fisica":
    count l in lessons where l.teacher == t
        and l.classroom.type == "lab_fisica": l == 1
\end{lstlisting}

\item Tetto cattedra Inglese: max 18h/settimana:
\begin{lstlisting}
forall t in teachers where t.subject == "Inglese":
    count l in lessons where l.teacher == t: l <= 18
\end{lstlisting}

\item Rossi non lavora il sabato (HARD):
\begin{lstlisting}
forall l in lessons where l.teacher == Rossi
                       and l.day == 6:
    false
\end{lstlisting}

\item Rossi: al pi\`u 1 sesta ora/settimana:
\begin{lstlisting}
count l in lessons where l.teacher == Rossi
                      and l.hour == 13: l <= 1
\end{lstlisting}

\item Tutti A050: max 5 ore/giorno:
\begin{lstlisting}
forall t in teachers where t.group == "A050":
    forall d in days:
        count l in lessons where l.teacher == t
            and l.day == d.index: l <= 5
\end{lstlisting}

\item Coupling Rossi 3A $\Rightarrow$ Bianchi 3B:
\begin{lstlisting}
(exists l in lessons:
    l.teacher == Rossi and l.class == 3A)
=> (exists l in lessons:
    l.teacher == Bianchi and l.class == 3B)
\end{lstlisting}

\item Rossi mai in palestra:
\begin{lstlisting}
forall l in lessons where l.teacher == Rossi:
    l.classroom.type != "palestra"
\end{lstlisting}

\item Tutti i prof di mate: almeno una lezione di mattina:
\begin{lstlisting}
forall t in teachers where t.subject == "Matematica":
    exists l in lessons where l.teacher == t:
        l.hour <= 9
\end{lstlisting}
\end{enumerate}

\subsection{Vincoli classe (5 esempi)}

\begin{enumerate}
\item La 1A non lavora il sabato:
\begin{lstlisting}
forall l in lessons where l.class == 1A: l.day != 6
\end{lstlisting}

\item Le prime: max 5 ore/giorno:
\begin{lstlisting}
forall c in classes where c.year == 1:
    forall d in days:
        count l in lessons where l.class == c.name
            and l.day == d.index: l <= 5
\end{lstlisting}

\item Quarte/quinte: due ore CONSECUTIVE di mate:
\begin{lstlisting}
forall c in classes where c.year >= 4:
    exists s1 in slots: exists s2 in slots:
        consecutive(s1, s2)
        and lesson(class==c.name, day==s1.day,
                   hour==s1.hour, subject==Matematica)
        and lesson(class==c.name, day==s2.day,
                   hour==s2.hour, subject==Matematica)
\end{lstlisting}

\item Linguistico anno 1: niente sabato:
\begin{lstlisting}
forall l in lessons
    where l.class.curriculum == Linguistico
      and l.day == 6: false
\end{lstlisting}

\item 1A: religione solo nelle prime 4 ore:
\begin{lstlisting}
forall l in lessons where l.class == 1A
                       and l.subject == Religione:
    l.hour <= 11
\end{lstlisting}
\end{enumerate}

\subsection{Vincoli aula (5 esempi)}

\begin{enumerate}
\item Lab Fisica 1: una sola lezione per slot:
\begin{lstlisting}
forall s in slots:
    count l in lessons
        where l.classroom == "Lab Fisica 1"
          and l.day == s.day and l.hour == s.hour:
        l <= 1
\end{lstlisting}

\item Matematica solo in aule taggate ``matematica'':
\begin{lstlisting}
forall l in lessons where l.subject == Matematica:
    "matematica" in l.classroom.tags
\end{lstlisting}

\item Storia delle terze: aula con proiettore:
\begin{lstlisting}
forall l in lessons where l.subject == Storia
                       and l.class.year >= 3:
    "proiettore" in l.classroom.tags
\end{lstlisting}

\item La palestra ospita SOLO Scienze motorie:
\begin{lstlisting}
forall l in lessons
    where l.classroom.type == "palestra":
    l.subject == Scienzemotorie
\end{lstlisting}

\item Lab chimica: max 2 lezioni concorrenti:
\begin{lstlisting}
forall s in slots:
    count l in lessons
        where l.classroom.type == "lab_chimica"
          and l.day == s.day and l.hour == s.hour:
        l <= 2
\end{lstlisting}
\end{enumerate}

\subsection{Vincoli materia (3 esempi)}

\begin{enumerate}
\item Mate distribuita: max 2 ore/giorno per classe:
\begin{lstlisting}
forall c in classes:
    forall d in days:
        count l in lessons
            where l.class == c.name
              and l.subject == Matematica
              and l.day == d.index: l <= 2
\end{lstlisting}

\item Mate preferibilmente entro la 5a ora:
\begin{lstlisting}
forall l in lessons where l.subject == Matematica:
    l.hour <= 12
\end{lstlisting}

\item Educazione fisica solo in palestra:
\begin{lstlisting}
forall l in lessons
    where l.subject == Educazionefisica:
    l.classroom.type == "palestra"
\end{lstlisting}
\end{enumerate}

\subsection{Vincoli gruppo / studenti (3 esempi)}

\begin{enumerate}
\item Studenti BES devono essere in un gruppo Sostegno (uso di
\texttt{exists g in s.groups}):
\begin{lstlisting}
forall s in students where "BES" in s.tags:
    exists g in s.groups: g.name == "Sostegno"
\end{lstlisting}

\item Studenti con debito mate quarto $\to$ gruppo Recupero:
\begin{lstlisting}
forall s in students
    where "debito_matematica_4" in s.tags:
    exists g in s.groups:
        g.name == "Recupero Matematica"
\end{lstlisting}

\item Studenti BES: nessuna lezione il sabato:
\begin{lstlisting}
forall l in lessons:
    forall s in students
        where l.class == s.class
          and "BES" in s.tags:
        l.day != 6
\end{lstlisting}
\end{enumerate}

\subsection{Vincoli relazionali / coupling (4 esempi)}

\begin{enumerate}
\item Rossi in Aula 12: solo Matematica:
\begin{lstlisting}
forall l in lessons where l.teacher == Rossi
                       and l.classroom == "Aula 12":
    l.subject == Matematica
\end{lstlisting}

\item Se Rossi in 3A allora deve usare lab fisica almeno una
volta:
\begin{lstlisting}
(exists l in lessons:
    l.teacher == Rossi and l.class == 3A)
=>
(exists l2 in lessons:
    l2.teacher == Rossi and l2.class == 3A
    and l2.classroom.type == "lab_fisica")
\end{lstlisting}

\item Inglese alle prime: max 3 ore/sett per (docente,classe):
\begin{lstlisting}
forall t in teachers where t.subject == "Inglese":
    forall c in classes where c.year == 1:
        count l in lessons
            where l.teacher == t
              and l.class == c.name: l <= 3
\end{lstlisting}

\item Religione alle prime: mai in 1a ora:
\begin{lstlisting}
forall l in lessons where l.subject == Religione
                       and l.class.year == 1:
    l.hour != 8
\end{lstlisting}
\end{enumerate}

\section{Errori comuni}

\begin{itemize}
\item \texttt{Atteso COLON, trovato OP (=)} -- usa \texttt{==}
nei confronti dopo \texttt{where} (\texttt{=} singolo e' lecito,
ma confonde il parser quando appare in contesti ambigui).
\item \texttt{sorgente sconosciuta 'foo'} -- sorgente non in
\texttt{\_VALID\_SOURCES} e non e' un path raggiungibile via
env. Controlla l'ortografia.
\item \texttt{oltre 1000000 atomi: probabile esplosione
combinatoria} -- quantificatori troppo annidati. Sposta filtri
nel \texttt{where} pi\`u esterno per ridurre il set; oppure
splitta il vincolo in pi\`u vincoli salvati separatamente.
\end{itemize}

Vedere \texttt{docs/general\_dsl.md} per la guida completa con
sezioni 1--12 (galleria di 30+ esempi commentati,
troubleshooting esteso, reference tabellare di tutti gli
attributi per ciascuna entit\`a).

\section{Compilazione DSL $\to$ CP-SAT (Step 1)}
\label{sec:dsl-cpsat}

A partire da Step 1 del piano multi-day, il DSL ha un
\emph{compilatore} verso CP-SAT
(\texttt{engine/dsl\_to\_cpsat.py}, classe
\texttt{DSLConstraintCompiler}). Mentre l'evaluator post-hoc
guarda una soluzione gi\`a costruita per dichiararla feasible o
violata, il compilatore aggiunge i vincoli \emph{durante} la
costruzione del modello CP-SAT, in modo che il solver eviti
strutturalmente le violazioni.

Il compilatore \`e usato in tre punti del sistema:
\begin{itemize}
\item dal \texttt{MonolithicSolver} (cap.~\ref{ch:metodo_cpsat})
  quando il flag \texttt{via\_dsl=True} \`e attivo;
\item dai pricer di Branch-and-Price
  (cap.~\ref{ch:tecniche_avanzate}) per applicare gli stessi
  vincoli al sotto-problema;
\item da \texttt{PhaseBDaySolver} quando si invoca con
  \texttt{via\_dsl=True} per riprodurre la pipeline legacy
  affiancata dal DSL.
\end{itemize}

Tutte le sorgenti delle lezioni gi\`a note (\texttt{teachers},
\texttt{classes}, \texttt{slots}, ecc.) sono accessibili al
compilatore allo stesso modo. La differenza chiave \`e che
quando una variabile slot \`e ancora indecisa, il compilatore
emette il vincolo CP-SAT corrispondente al posto di valutare
una verit\`a; quando un'espressione \`e \emph{completamente
statica} (per esempio \texttt{consecutive(s1, s2)} con \texttt{s1}
e \texttt{s2} legati staticamente nel \texttt{forall}) la
valutazione viene eseguita a tempo di compilazione e si emette
un vincolo solo se serve.

\section{Predicati di slot: \texttt{consecutive}, \texttt{same\_day}}
\label{sec:dsl-slot-preds}

Step 3a aggiunge due predicati built-in che il compilatore
risolve a tempo di compilazione quando entrambi gli argomenti
sono slot legati staticamente:

\begin{itemize}
\item \texttt{same\_day(s1, s2)} \`e vero se i due slot cadono
  nello stesso giorno della settimana.
\item \texttt{consecutive(s1, s2)} \`e vero se i due slot cadono
  nello stesso giorno e differiscono di un'ora ($|h_1-h_2|=1$).
\end{itemize}

Sono i mattoni con cui si esprimono ``ore consecutive'',
``stacco fra plessi'', ``coppia di ore quando ce ne sono due
nel giorno'', ``compresenza che deve avvenire nello stesso
giorno''. Gli argomenti accettati sono o tuple letterali
\texttt{(d, h)} oppure path puntati come \texttt{l.slot} che
restituiscono una \texttt{(day, hour)}.

Esempio: una classe deve avere mate in tre ore consecutive in
\emph{qualche} giorno della settimana:
\begin{lstlisting}
exists s1 in slots: exists s2 in slots: exists s3 in slots:
    consecutive(s1, s2) and consecutive(s2, s3)
    and lesson(class==3B, day==s1.day,
               hour==s1.hour, subject==Matematica)
    and lesson(class==3B, day==s2.day,
               hour==s2.hour, subject==Matematica)
    and lesson(class==3B, day==s3.day,
               hour==s3.hour, subject==Matematica)
\end{lstlisting}

\section{Path \texttt{l.classroom.plesso}}
\label{sec:dsl-classroom-plesso}

Ancora in Step 3a, l'evaluator post-hoc e il compilatore
risolvono il path chained \texttt{l.classroom.plesso}: per ogni
lezione \texttt{l}, l'aula assegnata viene letta dalla mappa
\texttt{classroom\_for\_slot} fornita da Phase B / Phase
classroom-assignment, e il plesso viene risolto via la tabella
\texttt{plessi\_data}. Senza queste due strutture il path
ritorna \texttt{None} e il compilatore emette un diagnostico.

L'esempio canonico \`e una commuting rule fra plessi (vedi
sez.~\ref{sec:plessi}): ``nessun docente pu\`o avere lezione in
Plesso B nell'ora successiva a una lezione in Plesso A'':
\begin{lstlisting}
forall l1 in lessons where l1.classroom.plesso == 1:
    forall l2 in lessons where l2.classroom.plesso == 2
        and l2.teacher == l1.teacher
        and consecutive(l1.slot, l2.slot):
        false
\end{lstlisting}

\section{Vocabolario di pattern canonici (Step 3b)}
\label{sec:dsl-canonici}

Step 3b introduce un \emph{vocabolario} di chiamate di funzione
che il compilatore riconosce in modo speciale ed emette
direttamente come gruppi compatti di vincoli CP-SAT. Ognuna di
queste forme corrisponde a un pattern frequente che, scritto in
DSL puro a colpi di \texttt{forall}, produrrebbe centinaia o
migliaia di clausole. La forma canonica \`e zero-drift rispetto
all'encoding hardcoded del solver legacy.

\begin{tabular}{lp{6.7cm}}
\toprule
\textbf{Pragma} & \textbf{Equivalente legacy} \\
\midrule
\texttt{no\_holes\_class("3B")} &
  catena non-crescente di booleani che vieta i buchi
  nell'orario settimanale della 3B (legacy
  \texttt{add\_class\_no\_holes}) \\

\texttt{class\_present\_at\_hour("3B", 11)} &
  se la 3B ha una lezione in un giorno, lo slot all'ora 11
  deve essere occupato (legacy HARD-3 / ``uscita dopo le 12'') \\

\texttt{class\_day\_load\_in("3B", 0, 4, 5, 6)} &
  in un giorno la 3B ha 0 ore (giorno libero) o tra 4 e 6 ore
  (legacy HARD-1+HARD-2) \\

\texttt{teacher\_max\_per\_day("Rossi", 5)} &
  Rossi ha al massimo 5 ore in un giorno
  (legacy \texttt{MAX\_PROF\_HOURS\_PER\_DAY}) \\

\texttt{cattedra\_max\_per\_day("Rossi", "3B", "Matematica", 2)} &
  la cattedra (Rossi, 3B, Mate) ha al massimo 2 ore in un
  giorno (legacy \texttt{MAX\_PER\_DAY\_TRIPLE}) \\

\texttt{subject\_pair\_must("3B", "Scienzemotorie")} &
  quando 3B ha 2 ore di motoria nello stesso giorno, devono
  essere consecutive (legacy HARD-B / \texttt{add\_motorie\_pair}) \\

\texttt{subject\_pair\_exists("3B", "Matematica")} &
  quando 3B ha 2+ ore di mate nello stesso giorno, almeno una
  coppia di consecutive deve esistere (legacy HARD-A /
  \texttt{add\_math\_italian\_pair}) \\
\bottomrule
\end{tabular}

\medskip

Inoltre il pragma \texttt{no\_holes\_all\_classes()} e
\texttt{class\_present\_at\_hour\_all(11)} agiscono in batch su
tutte le classi nel \emph{slot scope} corrente, utile da
incollare in cima alla lista di vincoli senza dover enumerare i
nomi.

\section{Plesso commuting via DSL (Step 3c)}
\label{sec:dsl-plesso-commuting}

Step 3c collega le \texttt{PlessoCommutingRule} al DSL: la
helper \texttt{plesso\_commuting\_rule\_to\_dsl(rule)} produce
una lista di clausole DSL equivalenti alla regola legacy. La
property zero-drift \`e verificata nella suite di test:
applicare la regola via DSL produce esattamente le stesse
combinazioni proibite della pipeline hardcoded
\texttt{plessi\_constraints}.

\section{Seed legacy HARD via DSL (Step 3d-3e)}
\label{sec:dsl-seed-legacy}

Step 3d-3e introduce \texttt{seed\_implicit\_hardcoded(profs)}
che, dato il dizionario delle cattedre, restituisce una lista
di stringhe DSL che rappresentano TUTTI i vincoli HARD
hardcoded del solver legacy. Combinato con il flag
\texttt{via\_dsl=True} di \texttt{MonolithicSolver} e
\texttt{PhaseBDaySolver}, permette di costruire un modello
CP-SAT completo \emph{esclusivamente dal DSL}: nessun ramo di
codice hardcoded \`e attivo. Le invarianti che si garantiscono:

\begin{itemize}
\item zero-drift sintattico: le tuple di slot proibite/forzate
  emesse dal path DSL coincidono bit-per-bit con quelle del
  path legacy;
\item zero-drift semantico: le soluzioni feasible accettate
  dai due path sono identiche, verificato sui benchmark
  \texttt{small}/\texttt{medium}/\texttt{huge};
\item determinismo: l'ordinamento delle stringhe DSL prodotte
  da \texttt{seed\_implicit\_hardcoded} \`e stabile, quindi
  l'ordine di inserimento dei vincoli nel modello CP-SAT \`e
  riproducibile run-to-run.
\end{itemize}

Questo \`e il \emph{vehicle} per la migrazione progressiva: nei
prossimi step il path hardcoded verr\`a deprecato a favore del
path DSL, e l'esistenza di \texttt{seed\_implicit\_hardcoded}
permette di farlo senza perdere la copertura dei vincoli.

\section{Stato dei vincoli SOFT}
\label{sec:dsl-soft-status}

Il DSL accetta sia regole HARD che regole SOFT (con peso
opzionale). A livello di \emph{evaluator post-hoc} il SOFT \`e
pienamente implementato: le violazioni di clausole SOFT
contribuiscono allo score globale della soluzione e sono
visibili nelle metriche restituite da \texttt{run\_metrics}.

A livello di \emph{compilatore CP-SAT} (Step 1+) il SOFT \`e
attualmente accettato dal parser ma il backend di compilazione
non emette ancora la traduzione a costo soft (penalit\`a
positiva sull'obiettivo). Il diagnostico
\texttt{soft constraint X: SOFT cost not yet wired} \`e emesso
dal compilatore per ogni regola SOFT incontrata; la regola
viene archiviata nel modello come \emph{label informativo}, ma
il solver non aggiusta l'obiettivo. Nei prossimi step la SOFT
verr\`a integrata: l'utente legge ora la guida tenendo presente
che il SOFT in DSL \`e semanticamente definito ma il suo
\emph{enforcement} via CP-SAT \`e ancora in corso.

In pratica: usare HARD per qualunque cosa che non possa essere
violata; per le SOFT, usare il DSL come ``contratto
dichiarativo'' che oggi la pipeline esegue post-hoc e domani
eseguir\`a anche durante la compilazione.

\section{Diagramma dell'architettura DSL}
\label{sec:dsl-archi}

\begin{verbatim}
+--------------------------------------------------+
|  UI (modal "+ Nuovo vincolo DSL")                 |
|  webui/frontend/.../GeneralConstraintsCard.svelte |
+--------------------+-----------------------------+
                     |
                     v
+--------------------+-----------------------------+
|  POST /api/constraints/general                   |
|  webui/backend/routers/constraints.py            |
+--------------------+-----------------------------+
                     |
                     v
+--------------------+-----------------------------+
|  Parser + AST                                    |
|  webui/backend/utils/general_dsl.py              |
+--------------------+-----------------------------+
                     |
                     +---------+---------+
                     v         v         v
              post-hoc    compiler    SOFT score
              evaluator   to CP-SAT   (post-hoc)
              (cattura    (Step 1+)
               violazioni
               HARD)
                     |         |
                     v         v
              +--------------------+
              |   CP-SAT solver    |
              |   (Mono / Pricer / |
              |    PhaseB / RF /   |
              |    Meta)           |
              +--------------------+
\end{verbatim}

Stesso AST, due backend: l'evaluator legge la soluzione e
risponde ``feasible/violato''; il compilatore aggiunge vincoli
al modello prima della solve. La grammatica \`e unica e i nomi
delle sorgenti sono identici nei due backend.


\chapter{Tecniche di ottimizzazione avanzate}

Oltre alle metaeuristiche classiche (LNS / SA / TS / ILS) il
progetto include cinque tecniche aggiuntive. Vedere
\texttt{docs/optimization\_strategies.md} per la specifica
completa.

\section{Adaptive Large Neighborhood Search (ALNS)}

\texttt{experiments/alns.py}. Sei destroy operator
(random\_window, day\_cluster, worst\_fit\_day, teacher\_day,
classroom\_day, single\_class\_week) e tre repair operator
(cp\_sat\_window, greedy\_by\_soft, bfs\_fill\_back) selezionati
con roulette wheel sopra exponentially-decayed scores;
acceptance SA-like.

\section{Variable Neighborhood Search (VNS)}

\texttt{experiments/vns.py}. Cicla quattro vicinati di dimensione
crescente (1-swap, 2-swap, 3-chain, k-opt con $k\in[4,6]$). Ad
ogni miglioramento riparte dal vicinato 1; termina quando
un'intera passata 1$\to$4 non trova nulla.

\section{Hall's theorem pre-check}

\texttt{experiments/diagnostics/hall\_check.py}. Diagnostico
sincrono pre-Phase-A. Tre controlli:
\begin{enumerate}
  \item per (class, subject) coverage,
  \item subject supply vs demand,
  \item Hall sampling random-subset con sottoinsiemi esclusivi.
\end{enumerate}
Esposto in \emph{tre} punti UI: bottone inline in PhaseACard,
riga in AdvancedTechniquesCard, sezione del tab Statistiche.

\section{Column Generation (Dantzig-Wolfe)}

\texttt{experiments/column\_generation.py}. Master LP via
\texttt{scipy.linprog} HiGHS; sotto-problema per docente.
Skeleton single-iteration per scuole molto grandi ($>200$
classi). Default OFF nella pipeline.

\section{Lagrangian Relaxation}

\texttt{experiments/lagrangian.py}. Rilassa i ponti
inter-cluster e li dualizza con multipliers $\lambda_b$;
aggiornamento via subgradient ascent
$\lambda_{k+1} = \max(0, \lambda_k + \alpha_k g_k)$ con
$\alpha_k = \alpha_0 / (1+k)$. Per-iteration refinement via SA.

% =========== DSL generico ===========

\chapter{Diagnostica statistica}
\label{ch:diagnostica_statistica}

\grokfig{fig_diagnostica_statistica_grafico}{Istogramma su
quaderno a quadretti, con gaussiana tratteggiata sovrapposta e
$\pi$Tantum che indica il picco.}

Il tab \texttt{/diagnostics} aggrega tutte le analisi sul
modello e sulla soluzione attiva. Vedere
\texttt{docs/diagnostics.md}. Visualizzazioni con Apache
ECharts (lazy-loaded), tema \emph{pitantum} (palette indaco /
oro / avorio / terra-di-siena).

\section{Sensitivity Monte Carlo}
Genera $N$ perturbazioni feasible (sequenze di mosse atomiche
random accettate solo se HARD soddisfatto) e calcola statistiche
SOFT. Coefficiente di variazione $< 0.05$ $\Rightarrow$
soluzione vicina all'ottimo locale.

\section{Analisi bipartito}
Tre metriche networkx: densita' del grafo classe$\times$docente,
modularita' Newman-Girvan greedy, top-K betweenness centrality.

\section{Correlazioni e regressioni}
Tre regressioni statsmodels: OLS \texttt{load\_vs\_gaps},
OLS \texttt{subjects\_vs\_soft}, Logit
\texttt{saturation\_vs\_sixth}. Output con coefficienti,
p-values, $R^2$ (o pseudo-$R^2$) e interpretazione testuale
auto-generata.

\section{Distribuzioni}
Cinque distribuzioni (carichi docenti, classi per giorno,
materia$\times$slot heatmap, occupancy slot, KS-test, chi-quadro)
+ goodness-of-fit.

\section{Run telemetry}
Tabella \texttt{run\_telemetry} (alembic
\texttt{a848420325b3}). I moduli solver pushano sample via
\texttt{utils.telemetry.collector(run\_id, phase)}. Il run detail
page (\texttt{/runs/[id]}) mostra il grafico objective vs tempo
(multi-linea per phase) + statistiche per stage + export JSON.

% =========== Cap 7: API ===========

\chapter{API REST}

Tutti gli endpoint sono sotto \texttt{/api/}. Backend espone $\sim$118
route. Lista non esaustiva delle p\`iu rilevanti.

\section{Pattern CRUD}

Ogni risorsa CRUD ha la stessa forma:

\begin{center}
\small
\begin{tabular}{l l l}
\toprule
Verbo & Path & Descrizione \\
\midrule
GET    & \texttt{/api/<entity>}      & lista (q + sort)  \\
GET    & \texttt{/api/<entity>/\{id\}} & dettaglio       \\
POST   & \texttt{/api/<entity>}      & create           \\
PUT    & \texttt{/api/<entity>/\{id\}} & replace        \\
DELETE & \texttt{/api/<entity>/\{id\}} & delete         \\
\bottomrule
\end{tabular}
\end{center}

\section{Esempi curl}

\begin{lstlisting}[style=shell]
# stato dataset
curl http://127.0.0.1:8000/api/dataset/state

# import small con tutti i pool
curl -X POST -H "Content-Type: application/json" \
  -d '{"profile":"small","use_optimized":true,
       "import_curricula":true,"import_classrooms":true,
       "import_students":true}' \
  http://127.0.0.1:8000/api/dataset/import-profile

# vincolo logico HARD su un docente
curl -X POST -H "Content-Type: application/json" \
  -d '{"expression":"mai aula:LabFisica@mar","kind":"hard"}' \
  http://127.0.0.1:8000/api/teachers/15/logical-unavailabilities

# bulk dry-run su 3 classi
curl -X POST -H "Content-Type: application/json" \
  -d '{"entity_ids":[1,2,3],"action":"add_logical",
       "payload":{"expression":"lun8 AND lun9","is_hard":true,
                  "soft_penalty":100},
       "on_conflict":"dry_run"}' \
  http://127.0.0.1:8000/api/bulk/classes/dry-run

# coverage settimana
curl 'http://127.0.0.1:8000/api/coverage/week?week_start=2026-04-27'

# event lessons (monitor)
curl http://127.0.0.1:8000/api/monitor/event/1/lessons
\end{lstlisting}

\section{Riferimento gruppi di endpoint}

\begin{itemize}
  \item \textbf{Anagrafiche}: \texttt{/api/teachers},
        \texttt{/classes}, \texttt{/classrooms},
        \texttt{/subjects}, \texttt{/curricula},
        \texttt{/students}, \texttt{/groups}.
  \item \textbf{Cattedre}: \texttt{/api/assignments} +
        \texttt{teachers-for-subject} + \texttt{loads}.
  \item \textbf{Vincoli logici}:
        \texttt{/api/\{teachers,classes,classrooms\}/\{id\}/logical-unavailabilities},
        \texttt{/api/curricula/\{cid\}/logical-constraints}.
  \item \textbf{Bulk ops}:
        \texttt{/api/bulk/\{entity\}/dry-run|apply}.
  \item \textbf{Schedule}:
        \texttt{/api/schedule/by-class|by-teacher|by-room|by-slot},
        \texttt{move-lesson}, \texttt{move-preview},
        \texttt{lesson/\{lid\}/classroom}.
  \item \textbf{Coverage}: \texttt{/api/coverage/week|cell},
        \texttt{/api/absences}, \texttt{/api/substitutions}.
  \item \textbf{Optimization}:
        \texttt{/api/optimize/assignment|phase-b|lns|sa|ts|ils|full|classroom-assignment}.
  \item \textbf{Monitor}: \texttt{/api/monitor/events|summary|constraints|conflicts},
        \texttt{event/\{aid\}/lessons|lesson/\{lid\}}.
  \item \textbf{Import}: \texttt{/api/import/\{entity\}},
        \texttt{/api/import/\{entity\}/template}.
\end{itemize}

% =========== Cap 8: Estendere ===========

\chapter{Estendere il sistema}

\section{Aggiungere una nuova tabella + CRUD}

\begin{enumerate}
  \item Modello in \texttt{webui/backend/models.py}.
  \item Pydantic in \texttt{schemas.py} (\texttt{XxxIn}, \texttt{XxxOut}).
  \item Router in \texttt{routers/xxx.py}. Copia un router esistente
        come template.
  \item Includi il router in \texttt{main.py}:
        \texttt{app.include\_router(xxx.router)}.
  \item Adapter DSL in \texttt{utils/list\_query.py}.
  \item Frontend: nuova route
        \texttt{webui/frontend/src/routes/xxx/+page.svelte},
        riusando \emph{SortableQueryableList}.
  \item Voce in \texttt{+layout.svelte} per la nav.
\end{enumerate}

\section{Migrazione idempotente}

Aggiorna \texttt{db.py::\_apply\_lightweight\_migrations}. Pattern:

\begin{lstlisting}[style=py]
if insp.has_table("my_table") \
        and not has_column("my_table", "new_field"):
    conn.execute(text(
        "ALTER TABLE my_table ADD COLUMN new_field VARCHAR(64)"
    ))
\end{lstlisting}

\section{Aggiungere un nuovo tipo di vincolo}

\paragraph{Per-cella.} Copia il pattern di
\texttt{TeacherUnavailability}: unique (entity\_id, day, hour),
columns state/penalty/reason. Aggiorna \texttt{models.py},
\texttt{schemas.py}, router, \texttt{optimization.py::\_availability\_constraints},
\texttt{routers/monitor.py::\_build\_constraints} (per il tab
Vincoli), \texttt{routers/bulk.py::ENTITY\_UNAV\_MODEL} se vuoi
supportare bulk.

\paragraph{Logico.} Aggiungi \texttt{entity\_type} a
\texttt{LogicalUnavailability}; aggiorna
\texttt{routers/logical.py::ENTITY\_MODEL} e \texttt{ENTITY\_TYPE};
modifica \texttt{routers/monitor.py::\_build\_constraints}; aggiungi
il branch nel solver.

\section{Aggiungere un nuovo step di ottimizzazione}

\begin{lstlisting}[style=py]
def my_new_step(...):
    params = dict(...)
    run_id = create_run("my_step", "Nome leggibile", profile, params)

    def target(rid: int):
        # carica dati, lancia solver, persiste
        update_run(rid, progress=0.5)
        ...
        update_run(rid, solution_id=sid, obj_value=v, metrics=m)
        update_run(rid, progress=1.0)

    start_thread(run_id, target)
    return run_id
\end{lstlisting}

Esponi via \texttt{routers/optimize.py} con un endpoint POST.
Frontend in \texttt{/optimize} per il bottone di lancio + log
streaming via \emph{RunLogPanel} (SSE).

\section{Testing}

Non c'\`e una test suite formale. Per smoke-testare end-to-end:

\begin{itemize}
  \item Backend boot pulito: \texttt{uvicorn backend.main:app}.
  \item Live curl: vedere \S7.2.
  \item Frontend \texttt{vite build} deve passare clean.
  \item Migrazione idempotente: rilancia su DB esistente per
        confermare zero data loss.
\end{itemize}

% =========== Appendice ===========
\appendix


\chapter{Repository GitHub}

URL: \url{https://github.com/gborghi/timetable}.

Repository privato. La landing page contiene il README di alto
livello con quickstart, tour della UI, sintassi vincoli,
architettura, repository layout.


\chapter{Reference}

\begin{itemize}
  \item Pisinger, "Where are the hard knapsack problems?", 2005.
  \item Kohl, Karisch, Patriksson, "Very Large-Scale Neighborhood
        Search Techniques", 2002.
  \item Burke, Kingston, Pepper, "A standard data format for timetabling
        instances", 2008 (S1877050919318800).
  \item Ahuja, Magnanti, Orlin, "Network Flows", 1993.
\end{itemize}



\end{document}
